% sections/prepotential.tex

\section{The gravitational prepotential}
\label{sec:prepotential}

The cubic determinant $\det(X)$ on $\hthree$ plays a dual role in the
self-modeling framework: it appears as a factor in the experiential
density $\rho_J$ and as the gravitational prepotential determining all
matter-gravity couplings. We prove that this double duty is forced by
algebraic uniqueness.

\subsection{The cubic norm}

For $X \in \hthree$ as in Eq.~(\ref{eq:h3o-element}), the cubic norm
(determinant) is
\begin{equation}
\label{eq:det3}
\det(X) = \alpha\beta\gamma
  - \alpha |x_1|^2 - \beta |x_2|^2 - \gamma |x_3|^2
  + 2\,\re\bigl((x_1 x_2) x_3\bigr).
\end{equation}
This is the unique (up to scale) $\Ffour$-invariant cubic polynomial
on $\hthree$.

\begin{theorem}[Springer, 1962~\cite{Springer1962}]
\label{thm:springer}
$\dim \mathrm{Sym}^3(27^*)^{\Ffour} = 1$.
That is, the space of $\Ffour$-invariant cubic forms on $\hthree$ is
one-dimensional.
\end{theorem}

The $\Ffour$-invariance of $\det(X)$ is a classical consequence of
Springer's theorem. As a consistency check, we have verified it
computationally under three generating subgroups of $\Ffour$:
\begin{itemize}
\item $S_3$ (permutations of diagonal entries): 60 tests, maximum
  error $7.1 \times 10^{-15}$.
\item $G_2$ (octonion automorphisms): 210 tests using 14 Schafer
  derivation generators~\cite{Schafer1966}, maximum error
  $2.0 \times 10^{-13}$.
\item $\Spin{9}$ (Peirce representation): 360 tests using 36 grade-2
  generators, maximum error $2.2 \times 10^{-5}$ at $\epsilon = 10^{-6}$
  (larger than the $S_3$ and $G_2$ residuals because the 36 grade-2
  generators are only \emph{approximately} in $\mathfrak{spin}(9)
  \subset\mathfrak{f}_4$; exact $\Ffour$-invariance of $\det$ is
  guaranteed analytically by Theorem~\ref{thm:springer}, so this row
  is a numerical sanity check, not the basis of the claim).
\end{itemize}

\subsection{The $d_{IJK}$ tensor}

The cubic norm defines a totally symmetric tensor $d_{IJK}$ in the
Peirce basis via
\begin{equation}
\label{eq:dijk}
\det(X) = d_{IJK}\, h^I h^J h^K,
\end{equation}
where $h^I$ are coordinates in the 27-dimensional Peirce basis
($I = 1$ for $\Vone$; $I = 2, \ldots, 17$ for $\Vhalf$;
$I = 18, \ldots, 27$ for $\Vzero$). The index $I = 0$ is reserved
for the graviphoton (from the gravity multiplet, not from the Jordan
algebra); see Sec.~\ref{sec:identification}.

The tensor $d_{IJK}$ has 106 nonzero entries out of 3654 distinct
triples (97\% sparse), organized into exactly two nonzero Peirce
blocks:
\begin{equation}
\label{eq:dijk-blocks}
d_{IJK} \neq 0 \;\Longleftrightarrow\;
\begin{cases}
(I, J, K) \in (\Vone, \Vzero, \Vzero): & 10 \text{ entries}, \\
(I, J, K) \in (\Vhalf, \Vhalf, \Vzero): & 96 \text{ entries}.
\end{cases}
\end{equation}
All other Peirce block combinations are exactly zero (verified
exhaustively, not sampled). The two nonzero blocks have clear physical
interpretations:

\begin{itemize}
\item $(\Vone, \Vzero, \Vzero)$: the observer's coupling to the
  spacetime metric. The 10 entries reproduce the $\det_2$ bilinear form
  on $\Vzero$ exactly. This is the gravitational self-coupling.

\item $(\Vhalf, \Vhalf, \Vzero)$: matter coupled to spacetime. The 96
  entries split as 48 spacetime (coupling through the four $\htwo{C}_u$
  directions) plus 48 internal (coupling through the six $W$-directions).
  All 16 matter fields in $\Vhalf$ couple to all 4 spacetime directions:
  the coupling is universal.
\end{itemize}

\subsection{The double duty corollary}

The cubic norm $\det(X)$ appears in two contexts within the self-modeling
framework:
\begin{enumerate}
\item As a factor in the experiential density
  $\rho_J = \det(X) \cdot (\Tr(X^2) - \tfrac{1}{3})$, governing which
  states support self-modeling.

\item As the prepotential of the Gunaydin-Sierra-Townsend MESGT
  (Sec.~\ref{sec:identification}), governing gravitational dynamics.
\end{enumerate}

\begin{corollary}[Double duty]
\label{cor:double-duty}
Both roles of $\det(X)$ are forced by the same algebraic uniqueness
constraint. Specifically:
\begin{enumerate}[label=(\alph*)]
\item The GST framework requires a cubic prepotential $V$ on $\hthree$
  that is invariant under $\Aut(\hthree) = \Ffour$.
\item By Theorem~\ref{thm:springer}, $\dim\,\mathrm{Sym}^3(27^*)^{\Ffour}
  = 1$, so $V = c \cdot \det(X)$ for some constant $c$.
\item The experiential density $\rho_J$ also requires an
  $\Ffour$-invariant cubic factor, which is again $\det(X)$.
\end{enumerate}
Therefore both roles use $\det(X)$ because it is the only cubic
$\Ffour$-invariant available. The two roles enter with different physical
interpretations (scalar function vs.\ prepotential) but share the same
algebraic source. Neither is derived from the other.
\end{corollary}

The proof is non-circular: the GST framework appears only as a
\emph{hypothesis} (``\emph{if} a cubic prepotential is needed, it must
be $\det(X)$''), not as a derivation step. The existence of the GST
Lagrangian provides the hypothesis; the uniqueness of $\det(X)$ provides
the conclusion.

\subsection{Quantum numbers}

The 16-dimensional $\Vhalf$ subspace carries one generation of Standard
Model fermions. The quantum number assignments (electric charge $Q$,
hypercharge $Y$, isospin $J_{3L}$, $J_{3R}$, $B - L$) have been
verified to match Paper~7 as a multiset, including color multiplicities
(12 quarks + 4 leptons = 16 states). This confirms that the $d_{IJK}$
coupling structure is compatible with the matter content derived
independently from the Peirce decomposition and chirality chain.
